% 参考文献。注意：main.tex 中 \referencescn 已通过 thebibliography 自动生成
% “参考文献”一级标题（不带编号），并已 \addcontentsline 进入目录，
% 所以本文件**不要再写一级标题**，否则会出现两个“参考文献”标题。

\begin{thebibliography}{99}
\bibitem{cumcm2026b} 全国大学生数学建模竞赛组委会. 2026 年高教社杯全国大学生数学建模竞赛 B 题：
                    机器狗搜索与清除干扰源[Z]. 2026.
\bibitem{diskcover} Disk covering problem. Wikipedia.
                    \texttt{https://en.wikipedia.org/wiki/Disk\_covering\_problem}.
\bibitem{twopt} 2-opt. Wikipedia. \texttt{https://en.wikipedia.org/wiki/2-opt}.
\bibitem{cormen} Cormen T H, Leiserson C E, Rivest R L, et al. Introduction to Algorithms[M].
                 3rd ed. Cambridge: MIT Press, 2009. （最近邻与 2-opt 局部搜索的 TSP 启发式）
\bibitem{proj} 本题求解代码与全部数值复现（\texttt{cumcm-2026-robot-dog} 仓库）：
              \texttt{reports/RESULTS\_REPORT\_T4.md} 为结果权威记录，图表见
              \texttt{figures/}. 2026-09.
\end{thebibliography}